\subsection{Περιγραφή κύριου pipeline}


Στην συνάρτηση \texttt{runPipeline} περνάμε τις εξής παραμέτρους:


\begin{description}
    \item[userInputData\_Original] Από τα πειράματα που θα περάσουμε στο σύστημα των μηχανικών μοντέλων,το κομμάτι των \texttt{userInputData}.
    \item[timeSeriesData\_BIG\_Original] Από τα πειράματα που θα περάσουμε στο σύστημα των μηχανικών μοντέλων,το κομμάτι των \texttt{timeSeries}.
    \item[models] Παιρνάμε τα μοντέλα και τα hyperparameters που θέλουμε να κάνουμε tuning,
    ώστε να βρούμε το μοντέλο για κάθε αισθητήρα που δίνει τα καλύτερα αποτελέσματα.
    \item[\detokenize{column_to_take}] Η στήλη που θα εξετάσουμε,που για τις εκτελέσεις του κύριου pipeline θα είναι πάντα το \texttt{VOC smoothed},
    ακόμα και σε περιπτώσεις που η τιμή δεν αντιπροσωπεύει το denoised VOC,που θα αναφερθούμε αργότερα σε αυτό.
    \item[\detokenize{extra_params}] Εμπεριέχει τις παραμέτρους που θα εισάγουμε στην συνάρτηση \linebreak \texttt{\detokenize{get_Data_To_Be_Used}},για το περαιτέρω φιλτράρισμα των texttt{userInputData} και \linebreak \texttt{\detokenize{ timeSeriesData_BIG_Original}}.
    \item[\detokenize{apply_functions_to_X_data}] Η συνάρτηση για την διαχείριση outliers και gradient, με βάση την επανάληψη που είμαστε στην \texttt{loopTroughFunctions}.
    \item[\detokenize{type_of_scaling}] Την μέθοδο scaling που θα εφαρμόσουμε στα δεδομένα πριν την εισαγωγή τους στα μηχανικά μοντέλα.
    Το scaling είναι απαραίτητα για τα περισσότερα μοντέλα εκτός του Random Forest Regressor,όπως και για το Principal Component Analysis (PCA).
    \item[\detokenize{scale_at_row}] Αν θα εφαρμόσουμε το scaling στις στήλες ή τις γραμμές.Από προεπιλογή το εφαρμόζουμε στις γραμμές,παρόλο το συνηθισμένο είναι το scaling να εφαρμόζεται στις στήλες. 
\end{description}

Θα αναφέρουμε τις συνάρτησεις που εμπεριέχει το runPipeline και τι κάνουνε.

\subsubsection{\texttt{\detokenize{create_The_Arrays_For_The_Models}}}

Αποτελεί την συνάρτηση που παίρνει τα userInputData\_Original και timeSeriesData\_BIG\_Original και επιστρέφει τους πίνακες που θα χρησιμοποιήσουμε έτοιμους για να τα εισάγουμε στην συνάρτηση hyperparameters tuning.

\paragraph{\detokenize{get_Data_To_Be_Used}}
Ξεκινούμε με το να κόψουμε από κάθε πείραμα,από τα \texttt{\linebreak timeSeries},όλα τα δεδομένα που είναι πάνω από τα 15 δευτερόλεπτα πριν από την εισαγωγή πηγής,
δηλαδή για τα δευτερόλεπτα που είναι κάτω από τα -15 δευτερόλεπτα,όπως και πάνω από τα 249 δευτερόλεπτα μετά την εισαγωγή της πηγής.
Όλα τα δεδομένα που πληρούν και τις δύο προϋποθέσεις,αποκτούν μέγεθος στηλών ίσον με 265. 
Όποιο πείραμα δεν πληρεί μία από τις δύο προϋποθέσεις,διαγράφεται από \texttt{timeSeries} και \texttt{userInputData}.

Επιλέγουμε αν θα κρατήσουμε ή όχι τα πειράματα που έχουν ανοιχτεί πόρτα,όπως και αν δεχτούμε ή όχι πειράματα με θέσεις πηγή ρύπου εκτός των 16 θέσεων που ορίσουμε να διεξάγουμε πειράματα.
Αν διαλέξουμε να κρατήσουμε και δεδομένα άλλων θέσεων,θα πετάξουμε μόνο τα δεδομένα που δεν έχουμε ορίσει με κάποια πηγή.

Η επιλογή του να είναι τα δεδομένα μόνο από τις θέσεις αισθητήρων που ορίσαμε ότι θα μελετήσουμε έχουν γίνει σε προηγούμενη φάση μες στο notebook.
Οπότε,αν είχαμε και άλλες θέσεις αισθητήρων για το δωμάτιο αυτό στα δεδομένα,δεν θα τα αποκλείαμε.

\paragraph{\texttt{\detokenize{build_X_keys_And_Y_target_Arrays}}}

Παίρνουμε από κάθε πείραμα,από την \texttt{userInputData}, 
το κλειδί με το οποία θα πάρουμε τα δεδομένα του πειράματος από την \texttt{timeSeries},
την θέση της πηγής ρύπου και τις αποστάσεις των ευκλείδειων αποστάσεων των πηγών από τους αισθητήρες.

Δημιουργούμε ξεχωριστούς πίνακες για το κάθε στοιχείο 
και διαχωρίζουμε και τους 3 πίνακες με τον ίδιο τρόπο σε train και test πίνακες της numpy,δηλαδή τα \texttt{ndarray},
ώστε να μπορούμε να διαχωρίσουμε τα πειράματα που θα χρησιμοποιήσουμε για την εκπαίδευση των μηχανικών μοντέλων,σε σχέση με όσα 
χρησιμοποιηθούν για έλεγχο του μοντέλου.

Έτσι έχουμε δημιουργήσει τους πίνακες Y,που είναι οι πίνακες που περιέχουν τους τις μεταβλητές στόχους του μηχανικού μοντέλου,
ή απλώς στόχους,γιατί το μυαλό είναι ο στόχος, που στην περίπτωση μας,οι στόχοι είναι και οι ευκλείδειες αποστάσεις και οι θέσεις πηγών,
όπου ο πρώτος πίνακας στόχων είναι για τα επιμέρους μηχανικά μοντέλα και ο δεύτερος για το συνολικό σύστημα.

Σημαντική σημείωση είναι ότι ο διαχωρισμός πραγματοποιείται με την επιλογή stratify,οπότε είμαστε σίγουρη ότι κάθε διακριτή
τιμή των ευκλείδειων διαστάσεων ανά αισθητήρα θα εμπεριέχεται στον Y test,τον πίνακα ελέγχου μεταβλητής στόχου δηλαδή.

\paragraph{\texttt{\detokenize{build_X_train_And_X_test}}}

Οι πίνακες X,δηλαδή οι πίνακες με τα δεδομένα εισόδου X,τα δεδομένα δειγμάτων ή πίνακας σχεδιασμού,δημιουργούνται πλήρως σε αυτή την συνάρτηση.

Αφού διαχωρίσουμε όλα τα δεδομένα των τωρινών \texttt{timeSeries} ανά αισθητήρα,θα φτιάξουμε τα X για το train και το X test.
Θα φτιάξουμε για το καθένα X πίνακες \texttt{ndarray} μορφής N \times 265,με N τον αριθμό των κλειδιών του X train ή του X test αντίστοιχα.
Με βάση των πίνακα των κλειδιών κάθε X,γεμίζουμε με την ίδια σειρά κάθε γραμμή με όλη την χρονοσειρά του πειράματος που έχει αυτό το κλειδί.

\paragraph{scaleXrows,ScaleX και μία σημαντική υποσημείωση για τις συναρτήσεις που επεξεργάζονται τους πίνακες X}

Εδώ επιλέγουμε αν θα κάνουμε scaling στις γραμμές ή στις στήλες.

Κατά κανόνα,το αλλαγή κλίμακας των δεδομένων (scaling) γίνεται άνα στήλη,ώστε να φέρει τα δεδομένα των στηλών στο ίδιο εύρος τιμών,όπου όλα τα μοντέλα 
που χρησιμοποιούμε,εκτός Random Forest Regressor,το έχουν σαν δεδομένο.
Αλλά επειδή μπορούν να υπάρχουν πολύ μεγάλες διαφορές μεταξύ των πειραμάτων για την ίδια στήλη,και είναι χαρακτηριστικό των δεδομένων και όχι outliers,
είδαμε ότι η μεταμόρφωση των δεδομένων ανά στήλη έτεινε να συμπιέσει πολλά δεδομένα.

Παρατηρήθηκε ότι κάνοντας normalization άνα δείγμα,συγκεκριμένα με την συνάρτηση \texttt{normalize} της \texttt{Scikit\-Learn},όπου διαπράττει  
L2 normalization το κάθε δείγμα,όπου φέρνουμε την ευκλίδεια νόρμα ίση με το 0,τα αποτελέσματα βελτιώνονται αισθήτα.
Οπότε για αυτό δίνουμε την επιλογή στο να πάρουμε scalers που δεν εφαρμόζονται κανονικά σε στήλες,
όπως ο \texttt{StandardScaler} και ο \texttt{MinMaxScaler}.

Κάθε συνάρτηση που επεξεργάζεται τα X δεδομένα,παίρνει πάντα σαν τους δύο βασικούς τις παραμέτρους τα X train και X text,και τα επιστρέφει μαζί.
Ο λόγος που τα παίρνουμε ζευγάρι και δεν κάνουμε απλώς σε κάθε πίνακα την ανάλογη επεξεργασία είναι γιατί υποτίθεται ότι
δεν ξέρουμε τα δεδομένα ελέγχου,οπότε πολλές συναρτήσεις,όπως στην περίπτωση μας οι scalers βάση στηλών,προσαρμόζουν τα δεδομένα τους
με βάση τα δεδομένα εκπαίδευσης και με αυτά τα δεδομένα,επεξεργάζεται και τα δεδομένα ελέγχου.

\paragraph{\texttt{\detokenize{modify_function}}}

Εδώ πραγματοποιούνται όλες οι διαφορετικές συναρτήσεις που έχουμε περάσει \linebreak
από την \texttt{loopTroughFunctions}.
Αυτό μας επιτρέπει να δοκιμάσουμε διαφορετικές συναρτήσεις επεξεργασίας των X δεδομένων,για τις ίδιες συνθήκες παραμέτρων που εισάγαμε στην αρχή.

Δυστυχώς,κάναμε την συνάρτηση να επαναλαμβάνει μόνο τις συναρτήσεις σχετικά με την διαχείριση των Outliers 
και το αν θα πάρουμε τα αρχικά δεδομένα ή τις παραγώγους,τα gradient των δεδομένων.

Για Outliers,δίνουμε την επιλογή να αφαιρέσουμε τα outliers του συνόλου των δεδομένων ή τα outliers ανά στήλη.
Οπότε αρκετές από τις συναρτήσεις που θέλαμε να δοκιμάσουμε τις πραγματοποιήσαμε μέσω πέρασμα έξτρα παραμέτρων \linebreak
στην \texttt{loopTroughFunctions} και την \texttt{runPipeline}, 
αφού πρότερα είχαμε προσθέσει την ανάλογη συνάρτηση στην \linebreak\texttt{\detokenize{create_The_Arrays_For_The_Models}}.

Παρόλα αυτά,είχε σημασία η εκτέλεση και η εξέταση των παραπάνω συναρτήσεων,εξετάζοντας τες σε διαφορετικές συνθήκες,σε σχέση με άλλες πιο δοκιμαστικές προσπάθειες.
Με κάποιες μικρές αλλαγές,θα μπορούσε να δεχτεί μη προδεσμευμένο τύπο διαφορετικών συναρτήσεων,χωρίς να χρειάζεται να ξέρουμε τον τύπο των συναρτήσεων από πριν

Από δοκιμές του pipeline,φάνηκε ότι έβγαζε καλύτερα αποτελέσματα εάν κάνουμε το scaling πριν κόψουμε τις ακραίες τιμές ή ψάξουμε για την παράγωγο,οπότε το εφαρμόσαμε με αυτή την σειρά.


\subsubsection{\texttt{printPCA2components,printTSNE}}

Απεικονίζουμε το principal component analysis (PCA) σε δύο στήλες και T-distributed Stochastic Neighbor Embedding,ώστε 
να έχουμε μια πρώτη εικόνα κατά πόσο τα πειράματα με ίδια ή κοντινή ευκλείδεια απόσταση 
αν έρχονται κοντά μεταξύ τους ή δεν υπάρχει κάποια ξεκάθαρη ομαδοποίηση μεταξύ των ίδιων ή κοντινών ευκλείδειων αποστάσεων
σε σχέση με άλλες περιπτώσεις.Δεν έχει κάποιο λειτουργικό σκοπό η συνάρτηση αυτή.

\subsubsection{\texttt{\detokenize{trainModels_PrintResults_AndGetPredictedValues}}}

Εδώ είναι που βρίσκομε το καλύτερο δυνατό μοντέλο με τις ανάλογες παραμέτρους,
να εισάγουμε τα X test και Y test και να πάρουμε τις προβλέψεις των μοντέλων για τις τιμές ελέγχου
και απεικονίζουμε σε διαγραμμά την διαφορά πρόβλεψης μεταξύ πραγματικής και τιμής πρόβλεψης.

\paragraph{ \texttt{\detokenize{run_grid_search_find_optimal_model}}}

Χρησιμοποιούμε το \texttt{GridSearchCV} της \texttt{Scikit\-Learn},
όπου ψάχνει διεξοδικά όλες τις παραμέτρους και τα μοντέλα που τις δίνονται για να "κουρδίσει" (κάνει tuning)
τις παραμέτρους των μοντέλων (hyperparameters).

Εφαρμόζουμε διασταυρούμενη επικύρωση 5 τμημάτων (5-fold cross validation), που σημαίνει ότι για κάθε συνδυασμό παραμέτρων και μοντέλων,
από τα δεδομένα εκπαίδευσης η \texttt{GridSearchCV} χωρίζει τα δεδομένα σε 5 κομμάτια,και εκτελεί 5 επαναλήψεις κρατώντας ένα διαφορετικό κομμάτι των δεδομένων σαν δεδομένα ελέγχου και τα
άλλα τέσσερα τα χρησιμοποιεί σαν δεδομένα εκπαίδευσης.
Από τις 5 επαναλήψεις,παίρνουμε το μέσο όρο της μετρικής αξιολόγησης που έχουμε διαλέξει,στην περίπτωση μας το \texttt{\detokenize{neg_mean_squared_error}},
που είναι το MSE με αρνητικό πρόσημο, 
και επιλέγουμε τον καλύτερο συνδυασμό μοντέλων και παραμέτρων που βρήκαμε.

Δύο σημειώσεις που έχουμε να κάνουμε,είναι ότι πρώτον, ο αριθμός των τμημάτων,δεν χρειάζεται να είναι 5,αλλά αποτελεί την πιο προεπιλεγμένη και συνηθισμένη επιλογή 
για τον αριθμό των τιμών.
Ταιριάζει απόλυτα με τον αριθμό των δειγμάτων ανά διακριτή τιμή ευκλείδειας απόστασης,οπότε χρησιμοποιούμε αριθμό τμημάτων ίσον με 5. 

Δεύτερο,και σημαντικότερο,to "μοντέλο" που ελέγχει η \linebreak \texttt{GridSearchCV} δεν χρειάζεται να είναι ένα μηχανικό μοντέλο απαραίτητα.
Μπορεί να είναι ένα \texttt{pipeline} της \texttt{Scikit\-Learn}.
Ένα pipeline αποτελείται μία σειρά μεθόδων μετασχηματισμού των δεδομένων με ένα προαιρετικό μοντέλο πρόβλεψης,που είναι τα μηχανικά μοντέλα που μελετάμε.

Οι μετασχηματιστές από τα αντικείμενα της κλάσης \texttt{TransformerMixin},
που στην \texttt{Scikit\-Learn} ονομάζονται transformers,αποτελούν έτοιμες μεθόδους,
όπως διάφορους scalers και PCA,αλλά μπορούμε να δημιουργήσουμε και τις δικές μας transformers κλάσεις.

Χάρη στην δυνατότητα εισαγωγής ενός pipeline στην \texttt{GridSearchCV},μπορούμε να βρούμε σε πόσες στήλες θα συμπηκνώσουμε τα δεδομένα μας
μέσω της PCA.
Για να κρατήσουμε ένα καλό λόγο μεταξύ δειγμάτων και στηλών που εισάγουμε στο πείραμα,θα κρατήσουμε τον μέγιστο αριθμό των στηλών στα δέκα.

\begin{center}
\fbox{\parbox{0.8\linewidth}{\centering
[3,4,6,8,10,'passthrough']
}}
\end{center}

Με την 'passthrough',μπορούμε να δοκιμάσουμε την περίπτωση να μην συμπυκνώσουμε τις στήλες και να αφήσουμε όλες τις στήλες όπως είναι.
Αυτή την επιλογή την εντάξαμε επειδή τα decision trees,δηλαδή τα Random Forest και Gradient Boosting Regressors,δεν χρειάζονται απαραίτητα να εφαρμόσουμε dimensionality reduction.

Από τα καλύτερο pipeline ανά αισθητήρα,επιστρέφουμε το όνομα του μηχανικού μοντέλο,το σκορ που απέδωσε,τις παραμέτρους που διαλέχθηκαν και
το ίδιο το συνολικό pipeline αντικείμενο, ώστε να το εκτελέσουμε στην \texttt{\detokenize{getPredictedValues}}.

\paragraph{\texttt{\detokenize{getPredictedValues}}}

Εδώ χρησιμοποιούμε τα μοντέλα που βρήκαμε για να πάρουμε τις ευκλείδειες αποστάσεις που προβλέπουν με βάση τα δεδομένα ελέγχου,εισόδου και μεταβλητών στόχου.

\paragraph{\texttt{\detokenize{plantAllBaseSensorsData}}}

Η συνάρτηση αυτή χρησιμοποιεί την \texttt{\detokenize{plot_}} \linebreak \texttt{\detokenize{_pred_vs_actual}} ώστε να δούμε πόσο κοντά είναι η πρόβλεψη των πραγματικών τιμών στις προβλεπόμενες τιμές.
Το κρίνουμε με βάση την απόσταση των σημείων από την διαγώνια γραμμή,όπου αναπαριστά την περίπτωση της τέλειας πρόβλεψης,με R\textsuperscript{2}=1.
Εδώ παρουσιάζουμε και τις μετρικές αξιολόγησης του καθενός μοντέλου.

\subsubsection{\texttt{\detokenize{findIntersectionOfCircles}}}

Εδώ εφαρμόζουμε την μέθοδο τριαπλεύρησης,βρίσκουμε την τελική θέση πηγής και σχεδιάζουμε τις προβλεπόμενες θέσεις πηγών ρύπου στο δωμάτιο.
Επιστρέφουμε πίσω τις τελικές μετρήσεις αξιολόγησής του συνολικού συστήματος.

\paragraph{\texttt{\detokenize{trilateration_least_squares}}}

Η τριαπλεύρεση σε περιπτώσεις που είχαμε τέλειες προβλέψεις θα ήταν να βρούμε το σημείο που διασταυρώνονται οι τρεις κύκλου που θα δημιουργούσαν οι ευκλείδειες αποστάσεις της ακτίνας.
Το σύστημα τριαπλεύρεσης εύρεσης σημείου διασταύρωσης τριών σημείων εκφράζεται με το παρακάτω σύστημα \cite{trilateration}:

\begin{align}
(x-x_1)^2 + (y-y_1)^2 &= r_1^2 \\
(x-x_2)^2 + (y-y_2)^2 &= r_2^2 \\
(x-x_3)^2 + (y-y_3)^2 &= r_3^2
\end{align}

Με $x_n,y_n$ την θέση του αισθητήρα με $Id=n$,εκφρασμένη σε \texttt{x axis} και \texttt{y axis}, $r_n$ την ευκλίδεια απόσταση του αισθητήρα με $Id=n$ 
και $x,y$ το σημείο διασταύρωσης των τριών κύκλων.

Όμως,σε πολλές περιπτώσεις, αφού τα επιμέρους μοντέλα θα κάνουν αναλυθείς προβλέψεις,οι κύκλοι δεν θα έχουν  σημείο διασταύρωσης.
Οπότε για αυτό χρησιμοποιούμε την μέθοδο ελαχίστων τετραγώνων 
\texttt{\detokenize{least_squares}} της \texttt{SciPy}, που βρίσκει την λύση που ελαχιστοποιεί το σφάλμα του παραπάνω συστήματος.

Τα ελάχιστα τετράγωνα αυτή είναι που χρησιμοποιείται για την επίλυση παρόμοιων προβλημάτων \cite{leastSquaresTrilateration}.
Σαν ξεκίνημα,δίνουμε την μέση τιμή της \texttt{x axis} και \texttt{y axis} όλων των αισθητήρων.
Έτσι,εντοπίζουμε την τελική λύση των μηχανικών μοντέλων.

\paragraph{\texttt{\detokenize{plot_pred_vs_actual}}}

Χρησιμοποιούμε πάλι την συνάρτηση σχεδιασμού που χρησιμοποιήσαμε για τα βασικά μοντέλα.
Αυτή τη φορά, εισάγουμε τις θέσεις πηγή ρύπου που έχει δύο στήλες,
για να σχεδιάζουμε τις πραγματικές θέσης πηγή ρύπου όπως και τις προβλεπόμενες πάνω στο σχεδιάγραμμα του δωματίου.

Χρησιμοποιείται για να δούμε οπτικά πόσο μακριά είναι η πρόβλεψη από την αληθινή τιμή.
Το ότι σε αρκετές περιπτώσεις το δωμάτιο θα είναι μικρότερο από το συνολικό διάγραμμα δεν είναι παράβλεψη,αλλά δείχνει ότι 
σε κάποια από τα βασικά μοντέλα,υπάρχει λάθος πρόβλεψη.
Μπορούμε να πούμε ότι η πρόβλεψη είναι και πάλι μες στο δωμάτιο,στο πιο κοντινό πραγματικό σημείο,
αλλά αν το διορθώναμε αυτόματα θα είχαμε πλαστή εικόνα για το σύστημα των μηχανικών μοντέλων.

\subsubsection{Τελικά Σχόλια}

Αυτό είναι το σύνολο του κύριο pipeline για regression μοντέλα.
Πραγματοποιήθηκε παρόμοιο pipeline για τα classification,το οποίο όμως δεν θα περιγραφεί.


